\chapter{\texorpdfstring{The proof of Proposition~\ref{prop:criterion}}{The proof of Proposition 1.7}}
\label{chap:criterion}

This appendix proves that $\Om_n \cong \Ct{n}$ if and only if $c_1, \dots, c_n$
are~$2$-independent, for every $a \in \QQ$ and every~$n$ (Proposition~\ref{prop:criterion}), which
the paper quotes from~\cite{paper}. The direction ``$\Leftarrow$'' and the lemmas behind it are
Section~1 of~\cite{paper}, formalized for a rational parameter; the direction ``$\Rightarrow$''
without any hypothesis on~$a$ is new (Lemma~\ref{lem:irred-of-iso}). The general facts about
multiquadratic extensions, wreath products and even polynomials used here are in
Appendix~\ref{chap:auxiliary}.

\section{The tower of splitting fields}

\begin{facts}[Kummer tower]
  \label{facts:kummer}
  \uses{def:iterates, def:galois}
  \lean{QuadraticIterates.splittingField_succ_eq_sup_adjoin,
    QuadraticIterates.sq_mem_splittingField_of_mem_rootSet_succ,
    QuadraticIterates.splittingField_le_succ, QuadraticIterates.splittingField_mono,
    QuadraticIterates.relfinrank_succ_le,
    QuadraticIterates.relfinrank_succ_eq_finrank_adjoin,
    QuadraticIterates.mem_rootSet_iteratedPoly_succ,
    QuadraticIterates.mem_rootSet_succ_of_sq_eq_sub, QuadraticIterates.exists_sq_eq_sub}
  \leanok
  $K_{n+1} = K_n(\sqrt{\alpha - a} : \alpha \text{ a root of } f_n)$, so $K_{n+1}/K_n$ is a
  multiquadratic extension and $[K_{n+1} : K_n] \le 2^{2^n}$.
\end{facts}
\begin{proof}
  \leanok
  \uses{lem:multiquadratic}
  The roots of~$f_{n+1}$ are the~$\beta$ with~$\beta^2 + a$ a root of~$f_n$, i.e.\ the square
  roots of the~$\alpha - a$; a multiquadratic extension generated by~$m$ square roots has degree
  at most~$2^m$ (Lemma~\ref{lem:multiquadratic}).
\end{proof}

\begin{lemma}[Odoni's embedding; \cite{paper}, Lemma~1.4]
  \label{lem:odoni}
  \uses{def:galois, def:wreath, facts:kummer}
  \lean{QuadraticIterates.nonempty_mulEquiv_succ_iff,
    QuadraticIterates.exists_injective_monoidHom_wreathPower,
    QuadraticIterates.isPGroup_galoisGroup, QuadraticIterates.pow_two_pow_apply_root,
    QuadraticIterates.card_wreathPower_succ, QuadraticIterates.nonempty_mulEquiv_iff_finrank_eq}
  \leanok
  $\Om_n$ embeds into~$\Ct{n}$; hence $\Om_n \cong \Ct{n}$ iff $[K_n : \QQ] = 2^{2^n - 1}$, and
  \[ \Om_{n+1} \cong \Ct{n+1} \iff \Om_n \cong \Ct{n} \text{ and } [K_{n+1} : K_n] = 2^{2^n}. \]
\end{lemma}
\begin{proof}
  \leanok
  \uses{lem:sylowwreath}
  Every element of~$\Om_n$ has~$2$-power order (its~$2^m$-th power fixes the roots of~$f_m$, by
  induction on~$m$), so~$\Om_n$ is a~$2$-group acting faithfully on at most~$2^n$ roots and
  embeds into a Sylow~$2$-subgroup of~$S_{2^n}$, which is~$\Ct{n}$
  (Lemma~\ref{lem:sylowwreath}). The rest is the degree count
  $\#\Ct{n+1} / \#\Ct{n} = 2^{2^n}$ together with Facts~\ref{facts:kummer}.
\end{proof}

\begin{lemma}[Irreducibility; \cite{paper}, Lemma~1.2]
  \label{lem:irred}
  \uses{def:cseq, def:iterates}
  \lean{QuadraticIterates.irreducible_iteratedPoly_succ_of_not_isSquare,
    QuadraticIterates.irreducible_iteratedPoly_of_not_isSquare_cSeq,
    Polynomial.Monic.exists_mul_comp_neg_X_eq_of_not_irreducible,
    Irreducible.isUnit_or_associated_of_dvd_comp_of_comp_neg_X_eq,
    Irreducible.isUnit_or_associated_of_dvd_comp_of_associated_comp_neg_X,
    QuadraticIterates.iteratedPoly_succ_comp_neg_X}
  \leanok
  If~$f_n$ is irreducible over~$\QQ$ and~$c_{n+1}$ is not a square in~$\QQ$, then~$f_{n+1}$ is
  irreducible. Hence, if none of $c_1, \dots, c_n$ is a square, then~$f_n$ is irreducible.
\end{lemma}
\begin{proof}
  \leanok
  \uses{lem:evenpoly, lem:involution}
  $f_{n+1} = f_n(X^2 + a)$ is even and has no nontrivial even divisor, because such a divisor
  would be $g(X^2 + a)$ with~$g$ a nontrivial divisor of~$f_n$. So if~$h$ is a nontrivial divisor,
  so is $h(-X)$, coprime to~$h$, and pairing the irreducible factors gives
  $\pm f_{n+1} = g(X) \cdot g(-X)$; the sign is the one relating $f_{n+1}(0)$ to~$c_{n+1}$, so
  $c_{n+1} = g(0)^2$.
  \formalnote{The first statement is \lname{irreducible\_iteratedPoly\_succ\_of\_not\_isSquare},
  the second \lname{irreducible\_iteratedPoly\_of\_not\_isSquare\_cSeq};
  Lemma~\ref{lem:irred-of-iso} uses the first in contrapositive form.}
\end{proof}

\begin{lemma}[Squares in~$K_n$; \cite{paper}, Lemma~1.5]
  \label{lem:squares}
  \uses{def:cseq, def:galois, def:twoindep, def:wreath}
  \lean{QuadraticIterates.not_isSquare_algebraMap_iff_twoIndependent_snoc,
    QuadraticIterates.isSquare_algebraMap_iff_exists_mul_prod,
    QuadraticIterates.isSquare_algebraMap_iff_exists_sq_eq,
    QuadraticIterates.TwoIndependent.finrank_adjoin_range_eq_two_pow,
    QuadraticIterates.card_monoidHom_eq_two_pow,
    QuadraticIterates.isSquare_algebraMap_cSeq}
  \leanok
  If $\Om_n \cong \Ct{n}$ and $c_1, \dots, c_n$ are~$2$-independent, then a rational~$c$ is a
  non-square in~$K_n$ iff $c_1, \dots, c_n, c$ are~$2$-independent.
\end{lemma}
\begin{proof}
  \leanok
  \uses{facts:kummer, lem:descent, lem:multiquadratic, lem:wreathab}
  $c_{k+1} = \prod_\beta (\beta - a)$ over the roots of~$f_k$ with multiplicity is a product of
  squares in~$K_{k+1}$, so $c_1, \dots, c_n$ are squares in~$K_n$, and they span
  an~$n$-dimensional subspace of the kernel of $\QQ^\times/\QQ^{\times 2} \to K_n^\times/K_n^{\times
  2}$. That kernel has dimension at most $\dim \Hom(\Om_n, C_2) = n$ (Lemma~\ref{lem:wreathab}); so
  it is spanned by the~$c_i$, and the descent Lemma~\ref{lem:descent} gives the claim.
\end{proof}

\begin{lemma}[The degree criterion; \cite{paper}, Lemma~1.6]
  \label{lem:degree}
  \uses{def:galois, def:cseq, facts:kummer}
  \lean{QuadraticIterates.relfinrank_succ_eq_two_pow_iff, QuadraticIterates.relfinrank_succ_eq_pow,
    QuadraticIterates.rootShift, QuadraticIterates.coe_rootShift,
    QuadraticIterates.rootShift_ne_zero, QuadraticIterates.prod_rootShift_eq_cSeq,
    QuadraticIterates.prod_rootSet_sub_eq_cSeq, QuadraticIterates.prod_aroots_sub_eq_cSeq,
    QuadraticIterates.rootRelations_rootShift_eq_bot_iff,
    QuadraticIterates.exists_ringEquiv_rootShift_smul,
    QuadraticIterates.sub_ratCast_ne_zero_of_mem_rootSet,
    QuadraticIterates.cSeq_succ_ne_zero_of_irreducible,
    QuadraticIterates.card_rootSet_iteratedPoly}
  \leanok
  If~$f_n$ is irreducible and $c_{n+1} \ne 0$, then $[K_{n+1} : K_n] = 2^{2^n}$ iff~$c_{n+1}$ is
  not a square in~$K_n$.
\end{lemma}
\begin{proof}
  \leanok
  \uses{lem:descent, lem:invariant, lem:multiquadratic, lem:relations}
  $[K_{n+1} : K_n] = 2^{2^n - \dim V}$, where $V \subseteq \FF_2^{2^n}$ is the space of
  relations~$\varepsilon$ with $\prod_j (\alpha_j - a)^{\varepsilon_j}$ a square in~$K_n$.
The~$2$-group~$\Om_n$ permutes the coordinates of~$V$ transitively (as~$f_n$ is irreducible), so
  $V \ne 0$ iff~$V$ contains the all-ones vector (Lemma~\ref{lem:invariant}), i.e.\ iff
  $\prod_j (\alpha_j - a) = c_{n+1}$ is a square in~$K_n$.
\end{proof}

\section{The criterion for irreducible iterates}

\begin{lemma}[The criterion for irreducible iterates; \cite{paper}, Section~1]
  \label{lem:criterion-irr}
  \uses{def:galois, def:wreath, def:twoindep, def:cseq}
  \lean{QuadraticIterates.section1_a_iff_b_of_irreducible}
  \leanok
  If $f_0, \dots, f_n$ are irreducible over~$\QQ$, then $\Om_n \cong \Ct{n}$ if and only if
  $c_1, \dots, c_n$ are~$2$-independent.
\end{lemma}
\begin{proof}
  \leanok
  \uses{lem:odoni, lem:squares, lem:degree}
  Induction on~$n$: $\Om_{n+1} \cong \Ct{n+1}$ iff $\Om_n \cong \Ct{n}$ and
  $[K_{n+1} : K_n] = 2^{2^n}$ (Lemma~\ref{lem:odoni}) iff $c_1, \dots, c_n$ are~$2$-independent
  and~$c_{n+1}$ is not a square in~$K_n$ (induction and Lemma~\ref{lem:degree}; $c_{n+1} \ne 0$
  because~$f_{n+1}$ is irreducible, \lname{cSeq\_succ\_ne\_zero\_of\_irreducible}) iff
  $c_1, \dots, c_{n+1}$ are~$2$-independent (Lemma~\ref{lem:squares}).
\end{proof}

\begin{lemma}[Maximality forces irreducibility]
  \label{lem:irred-of-iso}
  \uses{def:galois, def:wreath}
  \lean{QuadraticIterates.irreducible_iteratedPoly_of_nonempty_mulEquiv,
    QuadraticIterates.irreducible_iteratedPoly_succ_of_relfinrank_eq,
    QuadraticIterates.cSeq_succ_ne_zero_of_relfinrank_eq,
    QuadraticIterates.splittingField_zero_left}
  \leanok
  If $\Om_n \cong \Ct{n}$, then~$f_k$ is irreducible over~$\QQ$ for every $k \le n$.
\end{lemma}
\begin{proof}
  \leanok
  \uses{lem:odoni, lem:irred, lem:degree}
  Induction on~$n$; $f_0 = X$ is irreducible. Let $\Om_{n+1} \cong \Ct{n+1}$. By
  Lemma~\ref{lem:odoni}, $\Om_n \cong \Ct{n}$ and $[K_{n+1} : K_n] = 2^{2^n}$, and by induction
  $f_0, \dots, f_n$ are irreducible. Suppose~$f_{n+1}$ is reducible. Then~$c_{n+1}$ is a square
  in~$\QQ$ (Lemma~\ref{lem:irred}), hence in~$K_n$. If $c_{n+1} \ne 0$, Lemma~\ref{lem:degree}
  gives $[K_{n+1} : K_n] < 2^{2^n}$, a contradiction. If $c_{n+1} = 0$, then~$a$ is a root
  of~$f_n$ (as $c_{n+1} = \pm f_n(a)$), and an irreducible polynomial with a rational root has
  degree~$1$; so $n = 0$ and $a = 0$. Then $f_1 = X^2$ and $K_1 = K_0 = \QQ$, so
  $[K_1 : K_0] = 1 \ne 2$, a contradiction again.
\end{proof}
